\documentclass[10pt]{article}
\usepackage[margin=20mm]{geometry}
\usepackage{amsmath,amssymb,mathtools,fontspec,unicode-math,xurl,hyperref}
\IfFontExistsTF{DejaVu Serif}{\setmainfont{DejaVu Serif}}{\setmainfont{Latin Modern Roman}}
\setmathfont{latinmodern-math.otf}
\setlength{\parindent}{0pt}\setlength{\parskip}{6pt}\sloppy
\title{Initial scientific review of
An analytic invariant of \(G_2\) manifolds}
\author{Paper identifier: M25-0050}
\date{}
\begin{document}
\maketitle
Abstract: The assigned publisher version of Crowley, Goette and Nordström, Inventiones mathematicae 239 (2025), 865–907, is reviewed. One Major concern identifies an invalid variational inference in the small-eigenvalue argument on the main gluing proof chain. A completely specified operator countermodel proves the failure of that inference. A coercivity argument repairs the relevant cutoff contradiction on the diagonal \(r=\ell\); the full two-parameter estimates are not certified. Independent exact calculations reproduce the five lattice examples and the characteristic-class and Maslov-angle arithmetic. No counterexample to the asserted \(G_2\) gluing formula or disconnected-moduli examples is established.\par
\section*{Identity and declared coverage}
The reviewed object is the 43-page publisher PDF of [CGN], DOI 10.1007/s00222-024-01310-z, published online on 23 January 2025. Printed page is PDF page plus 864. The assigned manifest binds this PDF and a publisher metadata file; no separate computational supplement or author repository was identified in that manifest. Its actual source hash appears in Appendix 5 (\#app:record).\par
All 43 pages were read as text; 29 selected pages were also inspected as rendered images. This is a claim-focused initial review, not a claim that every geometric construction or cited theorem has been independently proved. The definitions and mod-48 relation in Section 1, gluing and configuration conventions in Section 2, all five numerical lattice examples in Section 3, the finite-dimensional signature calculation in Section 4, and the selected operator constructions and small-spectrum arguments in Section 5 were checked. The author’s preprint [CGNv5], dated 3 October 2024, was consulted at the passages relevant to the finding; it does not replace the frozen publisher version.\par
The construction requires closed spin seven-manifolds and compatible spin structures, ACyl Calabi–Yau pieces, the specified cyclic actions, and a hyper-Kähler rotation compatible with the torus gluing. The angle satisfies \(0<\vartheta<\pi\), and \(\rho=\pi-2\vartheta\in(-\pi,\pi)\). Spectral symmetry of the pieces is used only when \(k_\pm\leq2\). The general gluing formula retains the piece contributions. The rectangular case in Corollary 3 is read with the ordinary twisted-connected-sum convention clarified on p. 877; the review does not extend it to arbitrary cyclic quotients merely because \(\vartheta=\pi/2\).\par
\section*{Problem 1: an eigenvalue exclusion is used as a coercive estimate}
\subsection*{Location, claim, and dependency}
Equation reference: eq:sourcequotient. Proposition 5.4, pp. 898–899 (PDF 34–35), excludes small eigenspinors on each half under strict APS boundary conditions. For the negative half, the allowed trace is \(F^+(A)\), with no component in \(\ker A\). Proposition 5.5, pp. 899–902 (PDF 35–38), constructs spinors \(\psi_\pm\) that vanish at the outer boundary. Equations (54)–(55) imply, for one half, \[
 \|\psi_-\|\geq\tfrac12,\qquad
 \frac{\|D_{M_-,\ell,r}\psi_-\|}{\|\psi_-\|}
 \leq \frac{C c}{r}+O(e^{-r/C}).\] The first paragraph of p. 902 then infers a small Dirichlet eigenvalue of the first-order Dirac operator, contradicting Proposition 5.4. This inference is invalid for the boundary realization being used.\par
The dependence is explicit: Proposition 5.5 provides the spectral gap used in Proposition 5.7 (p. 905) to exclude spectral flow in the deformation from the modified operator to the geometric one. The final proof of Theorem 1 (p. 906) invokes Proposition 5.7 and the Dirac gluing theorem. Thus the defect concerns a core analytic justification, rather than an unrelated estimate. It is Major under the assigned core-derivation standard. This classification does not assert that Theorem 1 is false.\par
\subsection*{The domain obstruction}
Write the boundary Green form as \(\omega(u,v)=\langle u,\gamma v\rangle\), where \(\gamma A=-A\gamma\). The article itself states on p. 897 that the symplectic annihilator of \(F^+(A)\) is \[F^+(A)\oplus\ker A.\] Here \(\dim_\mathbb R\ker A=4\) (Remark 5.2). Consequently the strict APS realization is symmetric but is not self-adjoint: the adjoint admits additional zero-mode boundary traces. Imposing full Dirichlet trace zero also does not give a self-adjoint first-order Dirac realization. Adding a Lagrangian half of \(\ker A\) gives a different, self-adjoint modified APS problem, as distinguished earlier in Sections 4.1 and 5.2.\par
A small quotient \(\|T u\|/\|u\|\) bounds a singular value of \(T\), or an eigenvalue of \(T^*T\) with its own domain. It need not produce an eigenvalue of \(T\) with the original boundary condition. The paper requires precisely this unjustified change of spectral problem.\par
\subsection*{Complete countermodel to the inference}
Fix \(L>0\) and set \[\mathcal H=L^2((0,L);\mathbb C),\qquad
 T=-i\frac{d}{dx},\qquad \operatorname{dom}T=H^1_0((0,L);\mathbb C).\] This is a densely defined closed symmetric first-order Dirac operator with full Dirichlet trace zero. Integration by parts gives \(\operatorname{dom}T^*=H^1((0,L);\mathbb C)\), strictly larger than \(\operatorname{dom}T\). For any complex \(\lambda\), \(Tf=\lambda f\) implies \(f(x)=a e^{i\lambda x}\); the condition \(f(0)=0\) forces \(a=0\). Thus \(T\) has no nonzero eigenspinor for any \(\lambda\).\par
Nevertheless the admissible function \(f(x)=\sin(\pi x/L)\) satisfies \[\|f\|^2=\frac L2,\qquad
 \|Tf\|^2=\frac{\pi^2}{2L},\qquad
 \frac{\|Tf\|}{\|f\|}=\frac{\pi}{L}.\] Taking \(L>\pi/\varepsilon\) gives a quotient below any prescribed \(\varepsilon>0\), while the alleged first-order eigenvalue still does not exist. The second-order operator \(T^*T\) has the familiar Dirichlet eigenvalue \(\pi^2/L^2\); it is a different conclusion. This countermodel verifies every assumption of the abstract variational inference. It is not presented as an extra-twisted connected sum or as a counterexample to Proposition 5.5’s geometric statement.\par
\subsection*{A second limitation of the printed estimate}
Equation reference: eq:weak. Even granting the scalar lower bound used in Proposition 5.4, its last display only yields an inequality of the form \[
 \frac{a}{(\ell+r)^2}
 \leq C\{\lambda^2+e^{-\ell/C}+e^{-r/C}\}.\] It does not prove the claimed uniform assertion for all sufficiently large \(\ell\) and all \(r\geq r_0\), with fixed \(r_0\) and constants independent of \(\ell\). With \(r\) fixed and \(\ell\to\infty\), \((\ell+r)^2 e^{-r/C}\) is unbounded. With \(\ell\) fixed and \(r\to\infty\), \((\ell+r)^2 e^{-\ell/C}\) is unbounded. This remains a problem even after restricting to \(r\geq\ell\). For example, with all displayed constants set to one, \(\ell=100\), \(r=10^{60}\) and \(\lambda=0\), inequality [eq:weak] (\#eq:weak) is compatible with zero: its left side is about \(10^{-120}\) and its right side exceeds \(e^{-100}\). For arbitrary fixed constants, sufficiently disparate parameters give the same non-implication.\par
This checks the logic of the estimate, not the existence of a zero eigenspinor on the actual manifold. Along \(r=\ell\), the exponential errors are absorbable. The proof of Proposition 5.7 explicitly chooses that diagonal. A missing square in the preceding gradient display and a choice of constants in the scalar inequality do not resolve the domain obstruction or the stated two-parameter quantifier.\par
\subsection*{Verified repair for the cutoff contradiction on \(r=\ell\)}
The needed replacement is a lower bound for all Dirichlet spinors on a half, not merely exclusion of eigenvectors. Here is a proof in the geometry already constructed in Sections 2 and 4.2. It supplies the diagonal statement used below.\par
A half consists of a fixed connected compact core and an end whose length \(L\) is comparable to \(\ell+r\), with metrics uniformly equivalent to a fixed product metric. The finite quotient does not alter these estimates. For a scalar function \(f\) vanishing at the outer end, the fundamental theorem of calculus on the cylinder gives \[\int_{[0,L]\times X}|f|^2
 \leq C L^2\int_{[0,L]\times X}|df|^2,\qquad
 \|f|_{\{0\}\times X}\|^2
 \leq C L\int_{[0,L]\times X}|df|^2.\] On the fixed core \(K\), the trace Poincaré inequality \[\|f\|_{L^2(K)}^2\leq C_K
 \bigl(\|df\|_{L^2(K)}^2+\|f|_{\partial K}\|^2\bigr)\] follows, for example, by compactness: a violating normalized sequence would converge to a nonzero constant with zero boundary trace. Combining the estimates, with \(L\geq1\), gives \(\|f\|\leq C L\|df\|\), with \(C\) independent of the neck length. Kato’s inequality therefore implies \[\|\psi\|\leq C L\|\nabla\psi\|
 \quad\text{for every spinor with zero outer boundary trace.}\]\par
Equation reference: eq:coercive, eq:diagonal. Let \(D_{\rm geom}\) denote the geometric Dirac operator for the gluing metric. Equations (38)–(39) give \(D=D_{\rm geom}+E_\ell\) with \(\|E_\ell\|_\infty\leq C e^{-a\ell}\). The scalar curvature is zero off the transition regions and bounded below by \(-C e^{-a\ell}\) there. For a zero-trace spinor there is no boundary term in the integrated Lichnerowicz identity. Hence \[\begin{aligned}
 \|D\psi\|^2
 &\geq \tfrac12\|D_{\rm geom}\psi\|^2-\|E_\ell\psi\|^2 \notag\\
 &\geq \left(\frac{c_0}{L^2}-C e^{-a\ell}\right)\|\psi\|^2.
 \end{aligned}\] For \(r=\ell\) and \(\ell\) sufficiently large, this proves \[
 \|D_{M_\pm,\ell,\ell}\psi\|
 \geq\frac{c_1}{\ell}\|\psi\|.\] There is no appeal to a nonexistent Dirichlet eigenvalue of \(D\).\par
The cutoff construction (50)–(55) can now be used directly. The coefficients \(b,e,f\) are controlled by Green’s identity and \(\sin\vartheta\ne0\), and \(|a|\leq C/\sqrt r\). For \(r=\ell\), its nonzero boundary modes decay exponentially and its zero-mode cutoff error is \(O(c/\ell)\). Equation [eq:sourcequotient] (\#eq:sourcequotient) then contradicts [eq:diagonal] (\#eq:diagonal) after choosing \(c\) small enough and \(\ell\) large. The kernel spinor supplied by (40) remains. This proves the diagonal version of the gap assertion in Proposition 5.5.\par
The proof above establishes exactly that repair. It does not establish the uniform claim for arbitrary \(r/\ell\), nor replace every nonresonance and limit argument in the full Dirac gluing theorem. Those stronger statements remain to be supplied or explicitly restricted by the authors.\par
\subsection*{Metric deformation and the actual size of the perturbation}
Proposition 5.7, p. 905, also calls \(D_{\ell,r,w}-D_{\ell,r,1}\) an operator of degree zero for all \(w\in[0,1]\). The construction on pp. 893–895 changes the metric for \(0\leq w<1/2\). Such a difference is generally first order: the square of its geometric Dirac principal symbol records \(-|\xi|_{g_w}^2\). Equality of symbols would force equality of the metrics. Identifying the spinor bundles does not remove this fact. For \(w\geq1/2\) the metric is fixed, so the degree-zero description does apply there.\par
This does not obstruct the diagonal repair. The actual estimate is \[\|(D_w-D_1)u\|
 \leq C e^{-a\ell}\bigl(\|\nabla u\|+\|u\|\bigr).\] Uniform metric bounds and the integrated Lichnerowicz formula give \(\|\nabla u\|\leq C\|u\|\) for \(u\in\ker D_w\): \(D_{{\rm geom},w}u=-E_wu\), the zeroth-order term is bounded, and the scalar curvature has a uniform lower bound. Thus \(\|D_1u\|\leq C e^{-a\ell}\|u\|\) on that kernel. If \(\dim\ker D_w>1\), choose \(u\) in it perpendicular to the known one-dimensional kernel of the self-adjoint closed-manifold operator \(D_1\). The repaired diagonal gap contradicts this estimate. The known spinor \(s_{\ell,w}\) supplies the opposite inequality \(\dim\ker D_w\geq1\). Possible changes of volume density under the bundle identification are uniformly equivalent, or can be removed unitarily, with the same first-order bound.\par
This proves the kernel-control mechanism on the diagonal with the appropriate relative estimate. Exponentially small local eta-variation densities integrated over volume \(O(\ell)\) still have exponential decay after decreasing the exponent. A complete independent reconstruction of every eta-gluing and arbitrary-length limit is not claimed.\par
\subsection*{Counterevidence, versions, and precise impact}
The intended conclusion is plausible and has a substantial verified repair in the range actually chosen in Proposition 5.7. The invariant construction in Section 1 and the finite-dimensional Maslov arithmetic do not depend on the invalid variational inference. The review finds no incorrect numerical lattice value and no \(G_2\) manifold contradicting Theorems 1, 4 or 5.\par
The same strict-APS statement, Dirichlet-eigenvalue inference and degree-zero description occur in the targeted passages of [CGNv5]. Limited title/DOI/correction searches did not identify an applicable erratum; this is not proof that none exists. Bunke’s original theorem [Bunke] distinguishes self-adjoint Lagrangian boundary realizations and imposes regularity assumptions on its simplified formula. Its invertible-kernel simplification cannot simply be substituted here, since the two boundary Lagrangians share \(s\). The article’s additional treatment of that common zero mode is essential.\par
\section*{Independent checks and positive evidence}
\subsection*{Normalization of the analytic invariant}
Using formal tangent Chern roots \(x_1,\ldots,x_4\), the four positive-spin weights can be chosen as \[\tfrac12(x_1+x_2+x_3+x_4),\
 \tfrac12(x_1+x_2-x_3-x_4),\
 \tfrac12(x_1-x_2+x_3-x_4),\
 \tfrac12(x_1-x_2-x_3+x_4).\] Their product, with the paper’s orientation convention, is \((p_1^2-4p_2+8e)/16\). Substitution of \(\widehat A_8=(7p_1^2-4p_2)/5760\) and \(L_8=(7p_2-p_1^2)/45\) verifies equation (3) exactly: \(2e(S^+)=48\widehat A_8+e-3L_8\). Combining the APS index, signature and zero-count formulas yields the coefficients \(-24(\eta+h)(D)+3\eta(B)\) in Theorem 1.2.\par
For a parallel spinor the Berezin-degree parity argument in Lemma 1.3 gives a vanishing Mathai–Quillen term. The Lichnerowicz identity, \(SM\cong\mathbb R\oplus TM\), and the Bochner identity give \(h(D)=1+b_1(M)\) on the closed Ricci-flat manifold. Thus \(\nu=\bar\nu+24(1+b_1)\) modulo 48; the full-holonomy case has \(b_1=0\). These deductions verify integrality and local constancy within the stated holonomy locus, not on the space of arbitrary metrics.\par
The corrected homogeneous coefficients already acknowledged in Example 1.8 were not relabeled as a new issue. The rational contribution is exactly \[24\,\frac{12923}{2\cdot3^2\cdot5^6}
 -3\,\frac{4817}{3^2\cdot5^6}=1,\] and the displayed Chern–Simons terms cancel.\par
\subsection*{Maslov angles, endpoints, and lattice examples}
The image of boundary cohomology is Lagrangian by the pair exact sequence and Poincaré–Lefschetz duality. Paired non-middle degrees have conjugate spectra, so their Maslov contributions cancel. On \(H^3(X)\), the product decomposition and Hodge signs give the angles \(\alpha_j^+-\rho\) and \(\alpha_j^-+\rho\) in (35). Reducing these to \((-\pi,\pi]\) and excluding the endpoint \(\pi\) gives the endpoint weight one and the interior weight two in Proposition 4.7. The case \(\rho=0\) must be treated separately. The positive-space angles then contribute \(-\operatorname{sign}\rho\) and give \[m_{H^3(X)}=-16\rho/\pi+m_\rho,\qquad
 m_{\ker A}=\rho/\pi.\] Consequently \(3(-16\rho/\pi+m_\rho)-24\rho/\pi
=-72\rho/\pi+3m_\rho\), with the signs and factor 72 as printed.\par
For each example, the review formed exact rational orthogonal projections from the printed Gram matrix: \[P_E=E(E^tGE)^{-1}E^tG,\qquad A_E=2P_E-I,\] and computed the characteristic polynomial and signature of the span of the two polarising lattices. The product is the identity on its orthogonal complement.\par
Example | signature of span | characteristic polynomial | \(m_\rho\) | \(\bar\nu\)\par
3.6 | \((2,0)\) | \(x^2+1\) | \(-1\) | \(-39\)\par
3.7 | \((2,2)\) | \((x^2+1)^2\) | \(0\) | \(-36\)\par
3.8 | \((2,1)\) | \((x+1)(x^2+1)\) | \(0\) | \(-36\)\par
3.11 | \((2,0)\) | \(x^2-x+1\) | \(-1\) | \(-51\)\par
3.12 | \((2,2)\) | \((x^2-x+1)^2\) | \(0\) | \(-48\)\par
In Example 3.8 the \(-1\) eigenvector \((0,-3,1)\) has Gram norm \(-16\), so the \(\pi\) angle is in the negative space. The values for the first three examples use \(\rho/\pi=1/2\); the last two use \(2/3\). Orientation reversal changes their signs. The endpoint conventions were additionally checked in 13,338 exact rational-angle cases. These finite checks supplement the branch calculation; they do not prove geometric realization of arbitrary lists of angles.\par
For the examples used in Theorems 4 and 5, the topology inputs were cross-checked against [N], Examples 8.1 and 8.17: the ranks are 97 and 109 and the divisor of \(p_1/2\) is two. The two entries in [CHNP], Table 3, row \(b=74\), and the three entries at \(b=86\), have divisor of \(p_1\) equal to four. The table uses \(b_3=b+23\), explaining the row labels. These checks corroborate the cited classification data; the underlying building-block and classification proofs were not reconstructed. The target article gives \(\nu=12\) for the \(-36\) example and \(\nu=24\) for the \(-48\) example, by its own normalization. An inconsistent \(\nu=36\) sentence in the earlier external Example 8.1 is not imported as evidence against the reviewed article.\par
\subsection*{Operator calculations and excluded local notation}
The self-adjoint finite-rank corrections (38)–(39) annihilate the specified spinors, and the cylinder modes (45) satisfy the eigenvalue equation. For instance, \(\mu=13\), \(\lambda=5\) and \(\sqrt{\mu^2-\lambda^2}=12\) give an exact check obeying \(2|\lambda|<\mu\). The boundary equations following (62) have determinant \[\cos^2(z\pi/2)+\sin^2\vartheta\,\sin^2(z\pi/2)>0\] in the \((w,y)\) block. They give \(w=y=0\), with the remaining two zero modes only at \(z=1\). The two endpoint branches approach zero from opposite signs, giving the stated spectral-flow count once the analytic exclusion of other modes is justified. That last justification is not replaced by the finite-dimensional determinant calculation.\par
There is a separate local notation issue in the proof of Proposition 4.6, p. 889. With the conventional pullback \(R=\kappa^*\) on differential forms, \(RB=-BR\) already follows from \(R d=dR\) and \(R\bar *=-\bar *R\). The extra \(c_v\) in the printed \(c_v\kappa^*\) does not give the stated anticommutator. In an oriented orthonormal frame \((dx_1,\ldots,dx_6,dv)\), formula (26) gives \[(c_vR c_{x_1}+c_{x_1}c_vR)(1)=2\,dx_1\wedge dv\ne0.\] This was checked on the exterior algebra, including all 3,136 Clifford basis identities. The natural pullback \(R\) is the direct repair: it anticommutes with the odd signature operator and the boundary complex structure, commutes with \(A\), and preserves the cohomological image \(L_B\) by naturality. It therefore preserves the modified APS domain and proves spectral symmetry. A coefficientwise reflection in a separately specified product trivialization would be a different meaning of \(\kappa^*\). Because the geometric assertion is restored by the existing orientation-reversing isometry and a direct notation repair, this is recorded as an excluded local issue, not an additional substantive finding or a counterexample to Proposition 4.6. The spinor lift in Proposition 5.3 is explicitly defined differently.\par
\section*{Conclusion and outstanding work}
Item | Judgment | Verified repair or limit\par
Problem 1 | Major invalid core variational inference | Dirichlet coercivity and the kernel comparison proved on \(r=\ell\); full arbitrary-length estimates remain open in this review.\par
Invariant and lattice arithmetic | No additional issue found | Characteristic-class identity, mod-48 relation, branch terms and five reported values checked.\par
Spectral-symmetry notation | Excluded local notation | Natural pullback supplies a complete symmetry proof.\par
Acceptance and adoption | Pending / not requested | No human or independent specialist acceptance is claimed.\par
Remaining work includes full ACyl Calabi–Yau construction and matching theory, the complete topology and smooth-classification dependencies, a uniform two-parameter replacement for Propositions 5.4–5.7, and every analytic detail of the eta limits for the general \(k_\pm>2\) case. No empirical experiment is needed to check the demonstrated domain error. Further finite lattice computations would not discharge these analytic obligations. The scope supports a proof-revision request; it does not certify whole-paper correctness or disprove the paper’s principal conclusions.\par
\section*{Supporting checks}
The frozen publisher PDF SHA-256 is\par
> ec9dbcbf4d6f034c8a57ac20fd915602c50d4a5379e6ef32a64f7b7b9a672827.\par
It performs exact rational matrix and polynomial operations; only the explicitly labeled exponential-scale diagnostics use floating point. It executes no author software. The saved `independent\_results.json` has SHA-256\par
> 2e14a53455e89dfb00e0f21ce6bec418e6c9392af2d0d2188e41c99da4683e49.\par
Its successful run reports five Gram cases, 13,338 Maslov cases, 3,136 Clifford identities, both separated-cylinder modes, the characteristic-class identity and the Berger remainder one. The complete program and results accompany the report; these checks are not an empirical reproduction or a replacement for the analytic proofs above.\par
Source pages actually viewed as images were PDF pages 1, 5–8, 11, 13, 16–22, 25–28, 31–32, 34–42. Extraction, rendering and image inspection are distinguished in the evidence ledger. Web screenshot requests for external references returned text placeholders through the available interface; they are not credited as visual inspections. The primary-reference notes record the exact passages actually read and the failed direct arXiv PDF fetch, subsequently resolved for passage comparison through the primary HTML version.\par
Current compiler, support-input, log, TeX/PDF and rendered-page hashes are recorded in `build\_record.json` and the shared support manifest. No author code was executed or compiler executable installed.\par
The controller’s Chat/6Pro UI observation belongs only to that controller. Human scientific adjudication remains pending; author adoption has not been requested.\par
All durable outputs are confined to the assigned task staging tree. The report does not authorize final archiving or any external contact.\par
9 D. Crowley, S. Goette and J. Nordström, An analytic invariant of \(G_2\) manifolds, Invent. Math. 239 (2025), 865–907. DOI: 10.1007/s00222-024-01310-z. Assigned publisher PDF.\par
D. Crowley, S. Goette and J. Nordström, An analytic invariant of \(G_2\) manifolds, arXiv:1505.02734v5, 3 October 2024. Primary HTML: https://arxiv.org/html/1505.02734v5. Targeted comparison of Sections 4.3 and 5.3–5.5.\par
U. Bunke, On the gluing problem for the \(\eta\)-invariant, J. Differential Geom. 41 (1995), 397–448. DOI: 10.4310/jdg/1214456222. Sections 1.3 and 1.8, Theorems 1.9 and 1.17 consulted. Primary article copy: https://webhomes.maths.ed.ac.uk/\textasciitilde{}v1ranick/papers/bunke.pdf.\par
J. Nordström, Extra-twisted connected sum \(G_2\)-manifolds, Ann. Global Anal. Geom. 64, article 2 (2023). DOI: 10.1007/s10455-023-09893-1. Examples 8.1 and 8.17 consulted for topology inputs.\par
A. Corti, M. Haskins, J. Nordström and T. Pacini, \(G_2\)-manifolds and associative submanifolds via semi-Fano 3-folds, Duke Math. J. 164 (2015), 1971–2092. DOI: 10.1215/00127094-3120743. Bath-hosted accepted manuscript, Table 3 on manuscript p. 58 (PDF page 59, including repository cover), consulted.\par
\end{document}
